% ============================================================================
%  公共导言区：吉本斯《博弈论基础》习题答案详解
%  依赖：ctex + xelatex
% ============================================================================
\usepackage[a4paper,
            top=2.5cm, bottom=2.5cm,
            left=2.6cm, right=2.6cm,
            headsep=0.55cm, footskip=0.9cm]{geometry}

% ------------------------------ 数学 ------------------------------
\usepackage{amsmath, amssymb, amsfonts, mathrsfs}
\allowdisplaybreaks[4]
% 结论不使用黑色方框强调，保留公式内容本身。
\renewcommand{\boxed}[1]{#1}
% 正文段首统一缩进两个汉字；行间公式后的正文也恢复首行缩进。
\setlength{\parindent}{2em}
\makeatletter
\everydisplay{\aftergroup\@afterindenttrue}
\makeatother

% ------------------------------ 绘图与表格 ------------------------------
\usepackage{tikz}
\usetikzlibrary{matrix, positioning, calc, arrows.meta,
                decorations.pathreplacing, shapes.geometric}
\usepackage{array}
\usepackage{booktabs}
\usepackage{multirow}
\usepackage{makecell}
\usepackage{graphicx}

% ------------------------------ 版式 ------------------------------
\usepackage{titlesec}
\usepackage{fancyhdr}
\setlength{\headheight}{21.99997pt}
\usepackage[dvipsnames, table]{xcolor}
\usepackage{enumitem}
\usepackage{caption}
\usepackage{needspace}

% ============================================================================
%  颜色
% ============================================================================
\definecolor{maincolor}{HTML}{1F4E79}   % 主色：深靛蓝
\definecolor{accent}{HTML}{B03A2E}      % 强调色：砖红
\definecolor{rulegray}{HTML}{9AA5B1}    % 分隔线灰
\definecolor{lightbg}{HTML}{F5F7FA}     % 浅底
\definecolor{dotgray}{HTML}{3C3C3C}     % 节点黑

% ============================================================================
%  章节标题样式
% ============================================================================
\titleformat{\chapter}[display]
  {\normalfont\raggedright}
  {\sffamily\bfseries\large\color{maincolor}第\,\thechapter\,章}
  {7pt}
  {\sffamily\bfseries\Huge\color{maincolor!80!black}}
  [\vspace{7pt}{\color{maincolor}\titlerule[1.2pt]}]
\titlespacing*{\chapter}{0pt}{-10pt}{28pt}

\titleformat{\section}[hang]
  {\normalfont\sffamily\bfseries\large\color{maincolor}}{\thesection}{0.6em}{}
\titlespacing*{\section}{0pt}{1.6em}{0.7em}

% 页码与页眉
\pagestyle{fancy}
\fancyhf{}
\fancyhead[L]{\small\sffamily\color{maincolor!80!black}\nouppercase{\leftmark}}
\fancyhead[R]{\small\sffamily\color{maincolor!80!black}\thepage}
\renewcommand{\headrulewidth}{0.4pt}
\renewcommand{\headrule}{\hbox to\headwidth{\color{maincolor!45}\leaders\hrule height \headrulewidth\hfill}}
\fancypagestyle{plain}{\fancyhf{}\fancyfoot[C]{\small\sffamily\thepage}\renewcommand{\headrulewidth}{0pt}}

% ============================================================================
%  自定义命令
% ============================================================================

% ---- 习题标题：左侧色块 + 编号，自动进入目录 ----
\newcommand{\exer}[1]{%
  \Needspace{0.48\textheight}% 新题前至少保留约半页空间，避免题目标题落在页底
  \par\addvspace{1.5em}%
  \phantomsection
  \noindent
  {\color{accent}\rule[-0.18em]{4pt}{1.05em}}%
  \hspace{0.55em}{\sffamily\bfseries\large\color{accent}习题\,#1}%
  \par\nobreak
  \addvspace{0.45em}%
  \addcontentsline{toc}{subsection}{习题\,#1}%
  \ignorespaces
}

% ---- “解：”标记 ----
\newcommand{\sol}{\textbf{解：}}

% ---- “注：”标记（复核说明等） ----
\newcommand{\cmt}[1]{\par\smallskip\noindent{\small\textbf{注：}#1}\par\medskip}

% ---- 图：居中绘图 + 自动编号题注 ----
\newcounter{figcnt}[chapter]
\newcommand{\myfig}[2]{%
  \par\addvspace{0.8em}%
  \refstepcounter{figcnt}%
  \begin{center}%
    #1\\[5pt]%
    {\small\sffamily\color{rulegray}图\,\thechapter.\arabic{figcnt}\quad #2}%
  \end{center}%
  \addvspace{0.4em}%
}

% ---- 支付矩阵：统一风格的表格 ----
\newenvironment{paytab}{%
  \par\addvspace{0.6em}%
  \begin{center}\small
  \renewcommand{\arraystretch}{1.35}%
  \setlength{\tabcolsep}{6pt}%
}{%
  \end{center}%
  \addvspace{0.4em}%
}
% 矩阵外框样式
\newcommand{\mapline}{\arrayrulecolor{rulegray!70}\hline\arrayrulecolor{black}}

% ---- 支付矩阵左上角说明文字 ----
\newcommand{\pthead}[1]{\sffamily\bfseries #1}
% 支付矩阵单元格使用固定高度的盒子，使文字和分式都在格内真正垂直居中。
\newcommand{\paybox}[2]{\parbox[c][3.2em][c]{#1}{\centering #2}}

% ---- tikz 统一样式 ----
\tikzset{
  pnode/.style   = {circle, fill=dotgray!90, inner sep=0pt, minimum size=2.7mm},
  tnode/.style   = {circle, fill=dotgray!90, inner sep=0pt, minimum size=2.4mm},
  elab/.style    = {font=\small, fill=white, inner sep=1.2pt},
  pay/.style     = {font=\small, align=center},
  nodlab/.style  = {font=\small},
  branch/.style  = {line width=0.5pt, color=black!85},
  infoset/.style = {densely dashed, color=black!60, line width=0.6pt},
  contarc/.style = {line width=0.6pt, color=black!75},
}

% ============================================================================
%  目录与超链接
% ============================================================================
\usepackage{hyperref}
\hypersetup{
  colorlinks = true,
  linkcolor  = maincolor,
  citecolor  = maincolor,
  urlcolor   = maincolor,
  bookmarksnumbered = true,
  pdftitle   = {吉本斯《博弈论基础》习题答案详解},
  pdfauthor  = {若光ovo and ChatGPT},
  pdfsubject = {联系邮箱：a15759219744@163.com},
}
\usepackage{bookmark}

\setcounter{tocdepth}{2}
\setcounter{secnumdepth}{2}

\renewcommand{\contentsname}{目\quad 录}
